% !TeX root = ../slides.tex
% 01.tex — Customized MDP.

\begin{frame}{Minesweeper board as a customized MDP}
  \begin{itemize}
    \item Grid of $H \times W$ cells, $N$ hidden mines; revealed safe cells
          show the count of adjacent mines ($0$–$8$). \textbf{Flags are not considered} in this environment.
    \item Agent's goal: uncover every safe cell without ever selecting a mine.
  \end{itemize}

  \centering
  \begin{tikzpicture}
    % Row 0 — all hidden
    \msc{0}{0}{mshidden}{}
    \msc{1}{0}{mshidden}{}
    \msc{2}{0}{mshidden}{}
    \msc{3}{0}{mshidden}{}
    % Row 1 — the clicked mine (red); one more mine hides at the right,
    % indistinguishable from an ordinary hidden cell
    \msc{0}{-1}{mshidden}{}
    \msc{1}{-1}{mshidden}{}
    \msc{2}{-1}{msmine}{\textcolor{white}{\faBomb}}
    \msc{3}{-1}{mshidden}{}
    % Row 2 — numbered border; the left cell hides a third mine
    \msc{0}{-2}{mshidden}{}
    \msc{1}{-2}{msrevealed}{\msnum{2}}
    \msc{2}{-2}{msrevealed}{\msnum{2}}
    \msc{3}{-2}{msrevealed}{\msnum{2}}
    % Row 3 — safely revealed
    \msc{0}{-3}{msrevealed}{\msnum{1}}
    \msc{1}{-3}{msrevealed}{\msnum{1}}
    \msc{2}{-3}{msrevealed}{}
    \msc{3}{-3}{msrevealed}{}
  \end{tikzpicture}

  \vspace{0.1cm}
  \mslegend{mshidden}{Hidden} \qquad
  \mslegend{msrevealed}{Revealed / number} \qquad
  \mslegend{msmine}{Mine that ended the episode}

  \vspace{0.1cm}
  \small A $4\times4$ example after a loss. Three mines sit on this board.
\end{frame}

\begin{frame}{MDP modelling: the action space}
  Every Python class extending the \texttt{gym.Env} base class must have the attributes \texttt{action\_space}
  and \texttt{observation\_space} used to redefine $A$ and $S$.

  \vspace{-0.15cm}
  \begin{itemize}
    \item \textbf{Action space} - \texttt{gym.spaces.Discrete}. The action space on a board depends on the board $H$ and $W$. 
    It's an integer number, between 0 and \texttt{(H $\times$ W) - 1}.
    \item  Inside the \texttt{step()} method the action
    is mapped to a board cell by doing: $i = a \operatorname{//} W,\quad j = a \bmod W$.
    \item We could have chosen a \texttt{gym.spaces.MultiDiscrete} to directly represent an action as a tensor containing directly $[x y]$ coordinates,
     but this representation adapts more to DQN, who produces as much as $Q(s,a)$ as many are the possible actions.  
  \end{itemize}
\end{frame}

\begin{frame}{MDP modelling: the state space}
  \textbf{State space} - \texttt{gym.spaces.Box(low=-2, high=8, shape=H$\times$W)}: The state space of the custom MDP 
  corresponds to what the agent actually \textit{sees} when playing.
  A cell in the \textit{player board} could be in the following states:
  \begin{itemize}
    \item $-2$: the cell is unrevealed
    \item $-1$: this is \textbf{not actually a possible state} for the agent. It's only the way in which
    a mine is represented in the board containing the \textit{solution} of the game. The agent is informed
    about losing a game by the \textit{terminated} boolean variable of the \textit{Gym Env}.
    \item $0...8$: a revealed safe cell. The number represents the mines around the cell.
  \end{itemize}
\end{frame}

\begin{frame}{MDP modelling: the reward (first version)}
It's inspired by the paper \textcite{wang2025minesweepercnn}.
\begin{table}
    \centering
    \renewcommand{\arraystretch}{0.95}
    \begin{tabularx}{\textwidth}{l X r}
      \textbf{Outcome} & \textbf{Condition} & \textbf{Reward} \\
      \midrule
      Win & num. safe cells \texttt{==} num uncovered cells & $+ H \cdot W $ \\
      Loss           & mine selected                   & $- H \cdot W $ \\
      Progress       & informed safe cell uncovered     & $+1$ \\
      Guess          & safe cell \textit{discovered by chance}      & $-0.5$ \\
      Already open   & already-revealed cell selected   & $-0.5$ \\
    \end{tabularx}
  \end{table}
  However, during the learning phase, this design of reward presented \textbf{some issues}:
  \begin{enumerate}
    \item The agent learned to \textbf{remain} on an \textbf{already discovered} cell for many timesteps
    \item During \textit{a progression} more than one cell could
    be opened. However, the reward is still $+1$ \textbf{independently} on the number of the \textbf{opened cells}.
    \item Huge boards get huge rewards, little boards get little rewards.
  \end{enumerate}
\end{frame}

\begin{frame}{\textit{Aside}: Minesweeper's flood-fill}
  \twocol{
    \centering
    \textbf{Before} --- click $(2,2)$\\[0.15cm]

    \begin{tikzpicture}[xscale=0.45, yscale=0.45]
      \foreach \x in {0,...,4} {
        \foreach \y in {0,...,4} {
          \msc[minimum size=0.42cm]{\x}{-\y}{mshidden}{}
        }
      }

      \msc[minimum size=0.42cm,draw=badred,ultra thick]
           {2}{-2}{mshidden}{}
    \end{tikzpicture}

  }{
    \centering
    \textbf{After} --- cascade complete\\[0.15cm]

    \begin{tikzpicture}[xscale=0.45, yscale=0.45]

      % Row 0
      \msc[minimum size=0.42cm]{0}{0}{mshidden}{}
      \msc[minimum size=0.42cm]{1}{0}{mshidden}{}
      \msc[minimum size=0.42cm]{2}{0}{mshidden}{}
      \msc[minimum size=0.42cm]{3}{0}{mshidden}{}
      \msc[minimum size=0.42cm]{4}{0}{mshidden}{}

      % Row 1
      \msc[minimum size=0.42cm]{0}{-1}{mshidden}{}
      \msc[minimum size=0.42cm]{1}{-1}{msrevealed}{\msnum{1}}
      \msc[minimum size=0.42cm]{2}{-1}{msrevealed}{\msnum{1}}
      \msc[minimum size=0.42cm]{3}{-1}{msrevealed}{\msnum{1}}
      \msc[minimum size=0.42cm]{4}{-1}{mshidden}{}

      % Row 2
      \msc[minimum size=0.42cm]{0}{-2}{mshidden}{}
      \msc[minimum size=0.42cm]{1}{-2}{msrevealed}{\msnum{1}}
      \msc[minimum size=0.42cm]{2}{-2}{msrevealed}{}
      \msc[minimum size=0.42cm]{3}{-2}{msrevealed}{\msnum{1}}
      \msc[minimum size=0.42cm]{4}{-2}{mshidden}{}

      % Row 3
      \msc[minimum size=0.42cm]{0}{-3}{mshidden}{}
      \msc[minimum size=0.42cm]{1}{-3}{msrevealed}{\msnum{1}}
      \msc[minimum size=0.42cm]{2}{-3}{msrevealed}{\msnum{1}}
      \msc[minimum size=0.42cm]{3}{-3}{msrevealed}{\msnum{1}}
      \msc[minimum size=0.42cm]{4}{-3}{mshidden}{}

      % Row 4
      \msc[minimum size=0.42cm]{0}{-4}{mshidden}{}
      \msc[minimum size=0.42cm]{1}{-4}{mshidden}{}
      \msc[minimum size=0.42cm]{2}{-4}{mshidden}{}
      \msc[minimum size=0.42cm]{3}{-4}{mshidden}{}
      \msc[minimum size=0.42cm]{4}{-4}{mshidden}{}

    \end{tikzpicture}
  }
\vspace{0.2cm}
\begin{itemize}
  \item Imagine selecting a cell with no mines in its neighborhood (\textit{i.e. a 0 cell}).
  
  According to the game rules, all the adjacent connected cells are revealed.
  \item Instead of relying on an inefficent recursive solution, we though about a iterative solution
  \textit{stack-based}. The coords of the cells with no mine around are pushed
  as soon as they are found, and popped as soon as their boundaries are explored. No more pushes
  are done with one of this conditions: a mine is found, a numbered cell is found, the cell has already
  been uncovered by a previous visit.
\end{itemize}
\end{frame}

\begin{frame}{MDP modelling: action masking + reward normalization}
  \twocol{
    \textbf{Action masking}
      
      To let the agent \textit{risk} more, Q-values/logits
      are set to $-\infty$ before letting the agent perform the
      next action.
  }{
    \textbf{Reward normalization}

    With normalization, the rewards bounds are kept under control
    whatever is the size of the board game. 
  }

  \[
    R_t =
    \begin{cases}
    -1
    & \text{a mine is selected}
    \\[0.1cm]

    -\dfrac{0.5}{S}
    & \text{already unrevealed cell is selected}
    \\[0.18cm]

    \dfrac{\Delta n_t}{S}
    & \text{uncovering } {\Delta n_t} \text{ cells } \textbf{without guessing}
    \\[0.18cm]

    \dfrac{\Delta n_t}{S} - \alert{\dfrac{1.5}{S}}
    & \text{uncovering } {\Delta n_t} \text{ cells } \textbf{\alert{with a guess action}}
    \\[0.18cm]

    \textcolor{goodgreen}{1} + \dfrac{\Delta n_t}{S}
    & \text{\textcolor{goodgreen}{winning} } \textbf{without guessing}
    \\[0.18cm]

    \textcolor{goodgreen}{1} + \dfrac{\Delta n_t}{S} - \alert{\dfrac{1.5}{S}}
       & \text{\textcolor{goodgreen}{winning} } \textbf{\alert{with a guess action}}
    \end{cases}
  \]
\end{frame}